% Generated draft: requires mathematical and source review.
\documentclass[11pt]{article}
\usepackage[margin=25mm]{geometry}
\usepackage{fontspec}
\setmainfont{Latin Modern Roman}
\usepackage{amsmath,amssymb,mathtools,unicode-math}
\setmathfont{Latin Modern Math}
\usepackage{enumitem,needspace,xurl,hyperref}
\setlength{\parindent}{0pt}
\setlength{\parskip}{5pt}
\setlength{\emergencystretch}{3em}
\newcommand{\N}{\mathbb{N}}
\newcommand{\Z}{\mathbb{Z}}
\newcommand{\Q}{\mathbb{Q}}
\newcommand{\R}{\mathbb{R}}
\newcommand{\C}{\mathbb{C}}
\newcommand{\F}{\mathbb{F}}
\newcommand{\A}{\mathbb{A}}
\newcommand{\PP}{\mathbb{P}}
\newcommand{\E}{\mathbb{E}}
\newcommand{\eps}{\varepsilon}
\newcommand{\dd}{\,\mathrm{d}}
\begin{document}
\section*{UnsolvedMath ID: 11000156}
\textbf{Source problem code:} catalogue entry
\subsection{Definitions and mathematical statement}
Fix the source's group \(\operatorname{Map}_{g,n}\) and its designated boundary twist \(\delta\). Consider two ordered positive factorizations
\[
 \mathbf t=(t_{a_1},\ldots,t_{a_r}),\qquad
 \mathbf t'=(t_{b_1},\ldots,t_{b_s}),\qquad
 \prod_i t_{a_i}=\prod_j t_{b_j}=\delta,
\]
where every \(a_i,b_j\) is a nonseparating simple closed curve. Let \(X_{\mathbf t},X_{\mathbf t'}\) be the total spaces of the corresponding Lefschetz fibrations in the source construction. Assume
\[
 \chi(X_{\mathbf t})=\chi(X_{\mathbf t'}),\qquad
 \operatorname{sign}(X_{\mathbf t})=\operatorname{sign}(X_{\mathbf t'}).
\]
The question is whether a finite sequence of Hurwitz moves and the intended partial conjugations always transforms \(\mathbf t\) into \(\mathbf t'\).

A Hurwitz move may be written \((a,b)\mapsto(b,b^{-1}ab)\), with inverse allowed. In a standard product-preserving partial-conjugation convention, a consecutive block with product \(P\) is conjugated factorwise by \(h\) satisfying \(hPh^{-1}=P\), leaving the other factors fixed. The source does not define its partial-conjugation convention or restrictions on \(h\); this description is conditional, not an enlargement of the permitted moves. No separate global-conjugation or stabilization move is added.

The meaning of \(n\), the boundary-fixing convention, whether \(\delta\) is one twist or a boundary multitwist, and the precise fibration/section or pencil construction of the total spaces are missing. They must be fixed to make the two numerical hypotheses and move question unambiguous.
\subsection{Short English statement}
Do equal Euler characteristic and signature force the given nonseparating positive boundary-twist factorizations to be related by Hurwitz moves and the source's partial conjugations?
\subsection{Sources}
\begin{itemize}
\item \emph{Problems on Mapping Class Groups and Related Topics}, edited by Farb, Question 2.5, supplied PDF page 145. \url{https://www.math.uchicago.edu/~farb/papers/mcgbook.pdf}. Recorded access 2026-07-29; extraction pdf-font-structure. The preceding group, boundary-twist and partial-conjugation definitions are absent; no separate citation is supplied.
\item Local UnsolvedMath registry, \texttt{data/unsolvedmath/problems.json}, permanent ID \(11000156\), exact source code \texttt{catalogue entry}; supplied OPEN-TRIAGE summary says truncated.
\end{itemize}
\end{document}
